% Solution: Tadelis (2013) Exercise 11.3 "Even/Odd Symmetry". Input by \solitem. Check: check_T1.py (item 8).
\begin{exsol}{Tadelis Exercise 11.3}{236}{Foundations}{2}
\asked Solve the finite alternating-offer bargaining game of Section~11.2 when the number of periods $T$ is
even.

\method Backward induction with a common discount factor $\delta$ (Tadelis \S11.2, pp.~224--227).
Player~1 proposes in odd periods and player~2 in even periods, so for $T$ even \emph{player~2 makes the
last offer}. Let $s_m$ be the share of the pie, valued at the start of the subgame, that the proposer
obtains when $m$ periods remain.

\solution
\textbf{Last period.} The proposer takes the whole pie: $s_1 = 1$ (the responder accepts at indifference).

\textbf{Earlier periods.} With $m$ periods left, the responder gets $\delta s_{m-1}$ by rejecting, since next
period he is the proposer with $m-1$ periods left. The proposer offers him exactly that and keeps
\[
  s_m = 1 - \delta\, s_{m-1}, \qquad s_1 = 1 \quad\Longrightarrow\quad
  s_m = \frac{1 - (-\delta)^m}{1 + \delta}.
\]
(Induction: $s_m = 1 - \delta\,\frac{1 - (-\delta)^{m-1}}{1+\delta} = \frac{1 + \delta(-\delta)^{m-1}}{1+\delta}$ and
$\delta(-\delta)^{m-1} = -(-\delta)^m$.)

\textbf{$T$ even.} Player~1 is the proposer when $T$ periods remain. His offer is accepted in period~1
(agreement is immediate, as in Prop.~11.3, p.~226):
\[
  (v_1, v_2) = \Bigl(\frac{1 - \delta^T}{1 + \delta},\; \frac{\delta + \delta^T}{1 + \delta}\Bigr).
\]
For $T = 2$: $v_1 = 1 - \delta$ and $v_2 = \delta$. Strategies: with $m$ periods left the proposer offers the
responder $\delta s_{m-1}$, and the responder accepts iff his share is at least that.

\textbf{Comparison with $T$ odd} (the book's case, $v_1 = \frac{1 + \delta^T}{1 + \delta}$): for $T$ even the first
mover's share is \emph{below} $\frac{1}{1+\delta}$, and for $T$ odd it is above. Both tend to the
infinite-horizon share $\frac{1}{1+\delta}$ (\S11.3) as $T \to \infty$. Player~2, who has the last word when
$T$ is even, is ahead ($v_2 > \frac12$) iff $2\delta^T > 1 - \delta$. Example: $T = 2$ and $\delta = 0.8$ gives
$v_1 = 0.2$, $v_2 = 0.8$; with $\delta = 0.4$, $v_1 = 0.6$.

\interp The last word is worth a lot when patience is high ($\delta \to 1$): the last proposer can claim
almost the whole pie, so the earlier proposer must buy agreement with the value of that position. Being
first is an advantage only when delay is costly.

\verify \texttt{check\_T1.py} (item~8) checks the recursion against the closed form symbolically for
$T = 1,\dots,8$, the case $T = 2$, the equivalence ``player~2 ahead $\iff 2\delta^T > 1-\delta$'' on a grid,
and the limit $\frac1{1+\delta}$ with the ordering $\frac{1-\delta^4}{1+\delta} < \frac1{1+\delta} <
\frac{1+\delta^5}{1+\delta}$ (at $\delta = \frac12$). Tadelis prints no answer.

\mistakes
\begin{itemize}
  \item Using the odd-$T$ formula $\frac{1 + \delta^T}{1 + \delta}$. Parity decides who proposes last.
  \item Discounting with the wrong factor or from the wrong period: each earlier step uses $\delta s_{m-1}$,
  the responder's own continuation.
\end{itemize}

\answer{$s_m = \frac{1 - (-\delta)^m}{1 + \delta}$. For $T$ even, agreement in period~1 with
$(v_1, v_2) = \bigl(\frac{1 - \delta^T}{1 + \delta}, \frac{\delta + \delta^T}{1 + \delta}\bigr)$; e.g.\ $T = 2$ gives
$(1-\delta, \delta)$. The last mover (player~2) is ahead iff $2\delta^T > 1 - \delta$. The limit as $T \to
\infty$ is $\frac{1}{1+\delta}$.}
\end{exsol}
